% !PRINTED 222
\weberpage{222}

Wir wollen für den Augenblick zur Abkürzung
\weq{12}{
\varrho = \sqrt{-3 D}
}
setzen; dann folgt aus (7)
\weq{13}{
Q = 3 \varrho f(t),
}
und aus (11)
\weq{14}{
y = -\frac{1}{3} \sqrt[3]{3 \varrho f(t)},
}
und der Ausdruck (10) ergiebt
\weq{15}{
(3 a_0 x + a_1) f(t) = y f'(t) - 3 y^2.
}
Um die Uebereinstimmung dieses Resultats mit der \gesperrt{Cardanischen Formel} herzuleiten, gehen wir auf den \S.\,62 zurück und setzen in den dortigen Formeln $x = t$, $y = 1$.

Wegen (13) wird, nach \S.\,62, (11), (13)
\[
\eta = 0, \quad \xi = -h \varrho,
\]
und nach \S.\,62, (10)
\[
f(t) = \xi^3, \quad f'(t) = 3 \xi^2 \xi',
\]
worin $\xi' = h A_0$ nach \S.\,62, (13) die Ableitung von $\xi$ nach $t$ ist,
\[
f(t) = -h^3 \varrho^3, \quad f'(t) = 3 h^3 A_0 \varrho^2.
\]
Danach folgt aus (14)
\[
y = \frac{h \varrho \sqrt[3]{3 \varrho}}{3}
\]
und aus (15)
\[
3 a_0 x + a_1 = -h A_0 \sqrt[3]{3 \varrho} - \frac{1}{h \sqrt[3]{3 \varrho}},
\]
oder nach \S.\,62, (14), (15)
\weq{16}{
3 a_0 x + a_1 = A_0 (k - h) \sqrt[3]{3 \varrho},
}
worin nach \S.\,62, (15)
\[
h A_0 \sqrt[3]{3 \varrho} = \sqrt[3]{\frac{-q_0 + 3 \varrho a_0}{2}}
\]
\[
k A_0 \sqrt[3]{3 \varrho} = \sqrt[3]{\frac{+q_0 + 3 \varrho a_0}{2}}.
\]

Nehmen wir $a_0 = 1$, $a_1 = 0$ an, so erhalten wir aus (16) die \gesperrt{Cardanische Formel} (\S.\,35)
\weq{17}{
x = \sqrt[3]{\frac{-a_3}{2} + \sqrt{\frac{a_2^3}{27} + \frac{a_3^2}{4}}} + \sqrt[3]{\frac{-a_3}{2} - \sqrt{\frac{a_2^3}{27} + \frac{a_3^2}{4}}}.
}
